% !TEX program = xelatex
\documentclass[11pt,a4paper]{article}

\usepackage{fontspec}
\usepackage{polyglossia}
\setdefaultlanguage{english}

\usepackage{amsmath,amssymb}
\usepackage{graphicx}
\usepackage{booktabs}
\usepackage{array}
\usepackage{multirow}
\usepackage{xcolor}
\usepackage[colorlinks=true, linkcolor=blue, citecolor=blue, urlcolor=blue]{hyperref}
\usepackage{cleveref}
\usepackage[round]{natbib}
\usepackage{algorithm}
\usepackage{algpseudocode}
\usepackage{tikz}
\usetikzlibrary{shapes,arrows,positioning,fit,backgrounds}
\usepackage{subcaption}
\usepackage{listings}
\usepackage[normalem]{ulem}
\usepackage{siunitx}
\sisetup{group-separator={,}, group-minimum-digits=4}

\usepackage{bidi}

% --- Script Support ---

\setmainfont[
    UprightFont=LinLibertine_R,
    BoldFont=LinLibertine_RB,
    ItalicFont=LinLibertine_RI,
    BoldItalicFont=LinLibertine_RBI,
    Extension=.otf
]{LinLibertine}

\newfontfamily\arabicfont[
    Extension=.ttf,
    Script=Arabic,
    Scale=1.0,
    Renderer=Harfbuzz
]{Amiri-Regular}

\newfontfamily\bengalifont[
    Extension=.ttf,
    Script=Bengali,
    Scale=1.0,
    Renderer=Harfbuzz
]{NotoSerifBengali-Regular}

\newfontfamily\chinesefont[
    Extension=.ttc,
    FontIndex=0,
    Scale=1.0
]{NotoSerifCJK-Regular}

\newfontfamily\cyrillicfont[
    Extension=.otf,
    Scale=1.0
]{LinLibertine_R}

\newfontfamily\georgianfont[
    Extension=.ttf,
    Script=Georgian,
    Scale=1.0,
    Renderer=Harfbuzz
]{NotoSerifGeorgian-Regular}

\newfontfamily\greekfont[
    Extension=.otf,
    Scale=1.0
]{LinLibertine_R}

\newfontfamily\hebrewfont[
    Extension=.ttf,
    Script=Hebrew,
    Scale=1.0,
    Renderer=Harfbuzz
]{NotoSerifHebrew-Regular}

\newfontfamily\devanagarifont[
    Extension=.ttf,
    Script=Devanagari,
    Scale=1.0,
    Renderer=Harfbuzz
]{NotoSerifDevanagari-Regular}

\newfontfamily\koreanfont[
    Extension=.ttc,
    FontIndex=1,
    Scale=1.0
]{NotoSerifCJK-Regular}

\newfontfamily\malayalamfont[
    Extension=.ttf,
    Script=Malayalam,
    Scale=1.0,
    Renderer=Harfbuzz
]{NotoSerifMalayalam-Regular}

\newfontfamily\tamilfont[
    Extension=.ttf,
    Script=Tamil,
    Scale=1.0,
    Renderer=Harfbuzz
]{NotoSerifTamil-Regular}

\newfontfamily\telugufont[
    Extension=.ttf,
    Script=Telugu,
    Scale=1.0,
    Renderer=Harfbuzz
]{NotoSerifTelugu-Regular}

\newfontfamily\thaifont[
    Extension=.ttf,
    Script=Thai,
    Scale=1.0,
    Renderer=Harfbuzz
]{NotoSerifThai-Regular}

\newcommand{\textar}[1]{{\arabicfont\RL{#1}}}
\newcommand{\textbn}[1]{{\bengalifont #1}}
\newcommand{\textcn}[1]{{\chinesefont #1}}
\newcommand{\textjp}[1]{{\chinesefont #1}}
\newcommand{\textru}[1]{{\cyrillicfont #1}}
\newcommand{\textka}[1]{{\georgianfont #1}}
\newcommand{\textgeorgian}[1]{{\georgianfont #1}}
\newcommand{\textel}[1]{{\greekfont #1}}
\newcommand{\texthe}[1]{{\hebrewfont\RL{#1}}}
\newcommand{\texthebrew}[1]{{\hebrewfont\RL{#1}}}
\newcommand{\texthi}[1]{{\devanagarifont #1}}
\newcommand{\textdevanagari}[1]{{\devanagarifont #1}}
\newcommand{\textko}[1]{{\koreanfont #1}}
\newcommand{\textml}[1]{{\malayalamfont #1}}
\newcommand{\textta}[1]{{\tamilfont #1}}
\newcommand{\textte}[1]{{\telugufont #1}}
\newcommand{\textth}[1]{{\thaifont #1}}

\newcommand{\vecb}[1]{\mathbf{#1}}
\newcommand{\embdim}{128}

\title{Symphonym: Universal Phonetic Embeddings for Cross-Script Toponym Matching in Geospatial Data Integration}

\author{Stephen Gadd\\
\small School of Advanced Study, University of London, UK\\
\small Institute for Spatial History Innovation, University of Pittsburgh, USA\\
\small ORCiD: 0000-0003-3060-0181\\
\small \texttt{stephen.gadd@sas.ac.uk}\\
\small \texttt{stephen.gadd@pitt.edu}\\
\small \texttt{stephen@docuracy.co.uk (DO NOT PUBLISH)}}\\


\date{}

\begin{document}

\maketitle

\begin{abstract}
Integrating geographic information across languages and writing systems requires matching place names (toponyms) across different scripts. Existing approaches rely on language-specific phonetic algorithms or romanisation steps that discard phonetic information. This paper presents Symphonym, a neural embedding system that maps toponyms from 20 writing systems into a unified 128-dimensional phonetic space, enabling direct cross-script similarity comparison. A Teacher-Student architecture first learns from articulatory phonetic features derived from IPA transcriptions, then distils this knowledge into a character-level Student model requiring no phonetic resources or language identification at inference time. Trained on 32.7 million triplet samples from 67 million toponyms spanning GeoNames, Wikidata, and Getty TGN, the system achieves the highest Recall@1 (85.2\%) and MRR (90.8\%) on the MEHDIE cross-script benchmark, outperforming string baselines. An ablation using raw articulatory features alone achieves only 45.0\% MRR, confirming the neural training's contribution. Cross-script evaluation on 11,723 pairs across 170+ script combinations yields 90.7\% accuracy at the 0.75 similarity threshold. As a script-agnostic phonetic similarity layer, the system addresses a missing interoperability component in spatial data infrastructures. Deployed over 67 million toponyms with sub-second retrieval, the fixed-length embeddings enable operational geospatial data integration.
\end{abstract}

\textbf{Keywords:} toponym matching; phonetic embeddings; multilingual gazetteers; knowledge distillation; spatial data infrastructure

\noindent\textbf{Subject Classifications:} Cultural heritage; Geography; Geospatial data integration; GIS; Spatial data infrastructure


\section{Introduction}
\label{sec:introduction}

The vision of a comprehensive, queryable model of the planet's geography~\citep{goodchild2012digital,guo2020digital} depends on integrating place-referenced data from sources worldwide. National mapping agencies, OpenStreetMap, Wikidata, the Getty Thesaurus of Geographic Names (TGN), and historical archives all contain geographic references---but these span dozens of languages and more than 20 writing systems. The same city appears as ``London'', ``Лондон'', ``\textar{لندن}'', and ``\textcn{伦敦}'' across these sources. Linking such references is prerequisite to building integrated geospatial knowledge bases~\citep{hill2006georeferencing}, yet current computational methods cannot bridge script boundaries at scale.

The core difficulty is phonetic, not orthographic. Speakers recognise ``København'' and ``Copenhagen'' as the same city because the sounds are similar, not because the spellings match. Yet computational approaches to toponym matching have largely relied on string-based metrics (edit distance, Jaro-Winkler) or phonetic algorithms designed for single languages (Soundex for English, Cologne phonetic for German). These fail when names cross script boundaries: no edit distance metric can link ``\textcn{東京}'' to ``Tokyo''. Modern global geographic databases face this challenge directly: GeoNames alone contains 67 million toponyms in 20 scripts, and integrating them requires cross-script matching at scale.

Recent work has applied neural methods to geographic entity matching~\citep{qiu2024geobert,rama2016,mehdie2025}, but existing systems either operate within single scripts, require language identification at inference time, or target specific language pairs. No existing system places ``Νέο Μεξικό'' (Greek), ``\textbn{নিউ মেক্সিকো}'' (Bengali), ``\textar{نيومكسيكو}'' (Arabic), and ``Нью-Мексико'' (Cyrillic) near each other in embedding space using only raw character input. Early phonetic algorithms---Soundex~\citep{odell1918}, Metaphone~\citep{philips1990}, PHONIX~\citep{gadd1990phonix}---encode hand-crafted rules for specific Latin-script languages. While combining multiple string metrics improves within-script matching~\citep{recchia2013,santos2018}, no ensemble generalises across script boundaries. The cross-lingual embedding literature~\citep{conneau2018,artetxe2018} targets word meaning rather than phonetic form. \citet{mehdie2025} curate a benchmark of medieval Hebrew-Arabic toponym matches and build a specialist matching system for that language pair; their primary contribution is the benchmark itself---a carefully verified set of ground-truth matches curated by domain experts. Symphonym addresses a different scope: general-purpose cross-script matching across 20 writing systems without language-pair-specific resources.

The system presented in this paper is best understood not as a standalone matching algorithm but as an interoperability layer for multilingual geospatial infrastructure. Contemporary Spatial Data Infrastructures (SDIs) have achieved global standardisation of coordinate reference systems, spatial indexing, and metadata schemas, yet cross-script name reconciliation remains fragmented and ad hoc. Entity reconciliation in SDIs typically depends on shared identifiers, coordinate overlap, or string similarity within a common script; cross-script equivalence remains largely unresolved at scale. GeoNames, Wikidata, and TGN already function as nodes in the Linked Open Data cloud, but URI-based linkage presupposes that matching records have been identified---precisely the step that fails when names appear in different scripts. Embedding-based phonetic indexing introduces a reusable, language-agnostic bridging mechanism that complements identifier-based linkage, enabling more robust federation of distributed gazetteers.

The practical consequences of this gap are concrete. A researcher querying a consolidated gazetteer for ``Baghdad'' in Latin script cannot retrieve the Arabic \textar{بغداد}, the Cyrillic Багдад, or the Georgian \textgeorgian{ბაღდადი}---all phonetically near-identical renderings of the same name, but invisible to any string-matching algorithm because they share no characters. The same problem recurs at industrial scale: geographic databases worldwide record place names in their national scripts, and integrating these with Latin-script global gazetteers currently requires either manual curation or language-pair-specific transliteration pipelines that do not generalise. Historical geographic sources compound the difficulty: medieval travelogues, colonial-era surveys, and archival catalogues contain place names in scripts and orthographic conventions that differ from modern standard forms, and no existing SDI component bridges these systematically. Each of these scenarios---cross-script search, multi-source gazetteer federation, and historical source integration---requires the same underlying capability: a representation that captures phonetic similarity across script boundaries without language-specific engineering.

The scale of the system---67 million toponyms indexed with sub-second retrieval---is not merely an engineering achievement but a prerequisite for operational Digital Earth infrastructure. A phonetic similarity layer is useful only if it can be queried against the full scope of global gazetteer holdings without requiring batch pre-computation of pairwise scores. The combination of fixed-length embeddings, approximate nearest-neighbour indexing~\citep{malkov2020hnsw}, and a lightweight inference model ($<$1ms per query on CPU) meets this requirement, positioning the system as a candidate component in production SDI stacks.

This paper presents \emph{Symphonym}, a system that maps toponyms from any script into a unified phonetic embedding space. The key contributions are:

\begin{enumerate}
    \item A \textbf{script-aware but script-agnostic embedding architecture} mapping toponyms from 20+ writing systems into a unified \embdim-dimensional space where proximity reflects phonetic similarity.

    \item A \textbf{Teacher-Student training framework} that transfers phonetic knowledge from articulatory features derived from IPA transcriptions to a character-level model requiring no phonetic resources at inference time.

    \item \textbf{Phonetic similarity filtering} for training pairs using articulatory feature distance, preventing false equivalences between unrelated exonyms (e.g., ``Germany'' $\neq$ ``Deutschland'').

    \item A \textbf{three-phase curriculum} progressing from phonetic feature learning through knowledge distillation to hard negative discrimination.

    \item \textbf{Empirical validation} on the MEHDIE Hebrew-Arabic benchmark~\citep{mehdie2025}, achieving the highest Recall@1 (85.2\%) and MRR (90.8\%) of any tested method, and on 11,723 cross-script pairs spanning 170+ script combinations.

    \item \textbf{Production deployment} over 67 million toponyms with sub-second retrieval, demonstrating operational readiness for geospatial data integration workflows.
\end{enumerate}

The resulting system enables fuzzy phonetic search without language identification, script-specific preprocessing, or runtime phonetic conversion. Beyond cross-script matching, the approach naturally handles historical orthographic variation---the inconsistent spellings characteristic of pre-standardisation documents---and transfers effectively to personal names in archival sources, suggesting broad applicability to geospatial metadata integration where name forms vary across collections.

\paragraph{Scope and Application.} Symphonym addresses the \emph{name-matching} component of toponym resolution. The model operates purely on phonetic similarity between strings, with no access to geographic coordinates or spatial context. The contribution here is the phonetic similarity component alone. Within this layered architecture, phonetic similarity functions as a probabilistic prior rather than a final decision mechanism. Candidate sets retrieved via approximate nearest-neighbour search are subsequently intersected with spatial constraints (bounding boxes, administrative hierarchies), entity type filters, and temporal metadata. This staged design preserves high recall at the phonetic retrieval step while delegating precision to downstream geographic disambiguation---a division of labour that is critical for SDI workflows where recall losses at early stages are irrecoverable.

Concretely, the pipeline proceeds as follows: the 67 million toponyms from consolidated gazetteer sources are embedded by the Student encoder and pre-indexed via HNSW; at query time, an incoming toponym is encoded to the same 128-dimensional space ($<$1ms on CPU) and matched against the index; candidates are then filtered by geographic proximity, entity type, and temporal constraints to produce final matches.

In practice, the system will enable two important capabilities for users of the World Historical Gazetteer (WHG). First, researchers will be able to search for a toponym by entering an approximate phonetic rendering in their own language and script---a Greek scholar typing ``Ιεροσόλυμα'' will retrieve results for Jerusalem across Arabic, Hebrew, Latin, and Cyrillic attestations without needing to know the name in each script. Second, cultural heritage professionals using the WHG Reconciliation API will be better able to identify places mentioned in archival descriptions or transcripts of historical documents, where place names may appear in unfamiliar scripts or non-standard orthographies.


\section{Materials and Methods}
\label{sec:methods}

\subsection{Architecture Overview}

The system employs a Teacher-Student knowledge distillation architecture~\citep{hinton2015distilling} (Figure~\ref{fig:architecture}) where a Teacher network trained on articulatory phonetic features produces target embeddings that a Student network learns to approximate from character sequences alone. At inference time, only the Student is required, enabling deployment without phonetic resources.

Three principles guide the design: (1)~\textbf{Script awareness without script dependence}---the model handles 20 writing systems but produces embeddings in a unified space where script boundaries are transparent, achieved through deterministic script detection combined with learned script embeddings; (2)~\textbf{Phonetic grounding}---embedding similarity reflects phonetic rather than orthographic or semantic similarity, via the Teacher's articulatory feature space; (3)~\textbf{Inference efficiency}---the deployed model requires no runtime phonetic conversion, language identification, or external resources.

\begin{figure}[ht]
\centering
\begin{tikzpicture}[
    node distance=0.8cm and 1.2cm,
    box/.style={rectangle, draw, rounded corners, minimum width=2.2cm, minimum height=0.7cm, align=center, font=\small},
    widebox/.style={rectangle, draw, rounded corners, minimum width=3.2cm, minimum height=0.7cm, align=center, font=\small},
    input/.style={box, fill=blue!15},
    process/.style={box, fill=green!15},
    encoder/.style={widebox, fill=orange!20, minimum height=1.2cm},
    output/.style={box, fill=purple!15},
    loss/.style={box, fill=red!15, dashed},
    arrow/.style={->, >=stealth, thick},
    dashedarrow/.style={->, >=stealth, thick, dashed},
    label/.style={font=\footnotesize\bfseries},
]

% === TEACHER SIDE (left) ===
\node[input] (toponym1) {Toponym};
\node[process, below=of toponym1] (epitran) {G2P\\{\footnotesize (Epitran | Phonikud | CharsiuG2P)}};
\node[process, below=of epitran] (panphon) {PanPhon\\Features};
\node[encoder, below=of panphon] (teacher) {Teacher\\Encoder\\{\footnotesize BiLSTM + Attn}};
\node[output, below=of teacher] (temb) {$e_T \in \mathbb{R}^{128}$};

\draw[arrow] (toponym1) -- (epitran);
\draw[arrow] (epitran) -- (panphon);
\draw[arrow] (panphon) -- (teacher);
\draw[arrow] (teacher) -- (temb);

% === STUDENT SIDE (right) ===
\node[input, right=3.5cm of toponym1] (toponym2) {Toponym};
\node[process, below=of toponym2] (chars) {Characters\\+ Script};
\node[process, below=of chars] (lang) {Language\\(optional)};
\node[encoder, below=of lang] (student) {Student\\Encoder\\{\footnotesize BiLSTM + Attn}};
\node[output, below=of student] (semb) {$e_S \in \mathbb{R}^{128}$};

\draw[arrow] (toponym2) -- (chars);
\draw[arrow] (chars) -- (lang);
\draw[arrow] (lang) -- (student);
\draw[arrow] (student) -- (semb);

% === TRAINING PHASES (center) ===
\node[loss, below right=0.3cm and -0.5cm of temb] (phase1) {Phase 1:\\Triplet Loss};
\node[loss, right=0.8cm of phase1] (phase2) {Phase 2:\\Distillation};
\node[loss, right=0.8cm of phase2] (phase3) {Phase 3:\\Hard Neg.};

\draw[dashedarrow] (temb) -- (phase1);
\draw[dashedarrow] (temb) -- (phase2);
\draw[dashedarrow] (semb) -- (phase2);
\draw[dashedarrow] (semb) -- (phase3);

% === LABELS ===
\node[label, above=0.1cm of toponym1] {Training Path};
\node[label, above=0.1cm of toponym2] {Inference Path};

% === INFERENCE ANNOTATION ===
\node[font=\scriptsize, text=gray, right=0.2cm of student] (inf) {Deployed};
\draw[->, gray, thick] (inf) -- (student);

% === BACKGROUND BOXES ===
\begin{scope}[on background layer]
    \node[fit=(toponym1)(epitran)(panphon)(teacher)(temb),
          fill=blue!5, rounded corners, inner sep=0.3cm] {};
    \node[fit=(toponym2)(chars)(lang)(student)(semb),
          fill=green!5, rounded corners, inner sep=0.3cm] {};
\end{scope}

\end{tikzpicture}
\caption{Teacher-Student architecture. \textbf{Left}: The Teacher converts toponyms to IPA, extracts articulatory features, and encodes to 128-d embeddings (training only). \textbf{Right}: The Student processes raw characters with script/language metadata (deployed at inference). \textbf{Bottom}: Three training phases progressively transfer phonetic knowledge.}
\label{fig:architecture}
\end{figure}

\subsection{Teacher Network}
\label{subsec:teacher}

The Teacher encodes toponyms via IPA transcriptions and articulatory features (Figure~\ref{fig:architecture}, left). PanPhon~\citep{mortensen2016} represents each IPA segment as a 24-dimensional binary vector encoding universal articulatory properties (place, manner, voicing), independent of language identity. The phoneme /b/ receives identical features whether appearing in English ``Berlin'', Russian ``Берлин'', or Arabic ``\textar{برلين}''---all encode a voiced bilabial stop. This grounding in universal articulatory features enables generalisation: the Teacher learns relationships between articulatory configurations rather than specific languages, and this knowledge transfers automatically to any script representing the same sounds. The architecture comprises feature projection, bidirectional LSTM, multi-head self-attention, learned attention pooling, and output projection to \embdim{} dimensions with L2 normalisation.

\subsection{Student Network}
\label{subsec:student}

The Student processes raw character sequences with script and language metadata (Figure~\ref{fig:architecture}, right). Script is detected deterministically from Unicode code points, mapping characters to one of 20 script categories. The vocabulary contains \textbf{113,280 tokens} observed across 66.9 million toponyms. Each character maps to a 64-dimensional embedding; script and language contribute 16-dimensional embeddings each; and a \textbf{length bucket embedding} (one of 16 buckets encoding sequence length) provides an 8-dimensional conditioning signal broadcast across all positions. These are concatenated to form a 104-dimensional input per character.

\paragraph{Length-Aware Representation.} Toponym lengths vary dramatically---from two-character abbreviations to long institutional names---and naive similarity comparison across length disparities produces spurious matches. The Teacher's PanPhon192 representation (8 positional bins $\times$ 24 articulatory features) inherently produces sharper phonetic profiles for short names and more averaged profiles for long names, since the same fixed bin count must accommodate variable-length sequences. The Student's length bucket embedding directly addresses this: by conditioning every character representation on a discretised length signal, the model learns to calibrate similarity scores relative to string length, penalising matches between strings of very different lengths. This design is validated empirically: in production testing, the Student model substantially mitigates the length-sensitivity inherent in the Teacher's fixed-bin pooling.

During training, known languages are randomly replaced with \texttt{<UNK>} (50\% probability), forcing script-intrinsic pattern learning. The architecture mirrors the Teacher: BiLSTM, self-attention, attention pooling, and output projection with L2 normalisation. Character-level noise augmentation (insertions, deletions, substitutions, transpositions at 30\% probability) trains robustness to OCR errors and historical spelling variation.

\subsection{Training Curriculum}
\label{subsec:training}

Symphonym employs a three-phase curriculum.

\paragraph{Phase 1: Teacher Training.} The Teacher learns to produce embeddings where phonetically similar toponyms cluster together, using triplet margin loss:
\begin{equation}
\mathcal{L}_{\text{triplet}} = \max\!\big(0,\; \|e_T^a - e_T^p\|_2 - \|e_T^a - e_T^n\|_2 + m\big)
\label{eq:phase1-loss}
\end{equation}
where $a$, $p$, $n$ are anchor, positive, and negative, $m = 0.3$, and $\|\cdot\|_2$ is Euclidean distance. Script-aware negative sampling draws 80\% of negatives from the same writing system as the anchor, forcing fine-grained phonetic discrimination within scripts. Training uses AdamW ($\text{lr} = 10^{-4}$, weight decay $10^{-5}$) with cosine annealing over 50 epochs (33.5h on NVIDIA L40S GPU, final val\_loss 0.0056).

\paragraph{Phase 2: Student-Teacher Alignment.} The Student learns to approximate frozen Teacher embeddings, minimising a combined distillation loss:
\begin{equation}
\mathcal{L}_{\text{distill}} = \alpha \cdot \text{MSE}(e_S, e_T) + \beta \cdot \big(1 - \cos(e_S, e_T)\big)
\label{eq:phase2-loss}
\end{equation}
where $e_S$ and $e_T$ are Student and (detached) Teacher embeddings, and $\alpha = \beta = 1.0$. Language dropout and noise augmentation force script-intrinsic learning and input robustness. Training uses AdamW ($\text{lr} = 10^{-4}$) over 50 epochs (1.5h, final val\_loss 0.0591). Student-Teacher cosine similarity reached 0.942.

\paragraph{Phase 3: Discriminative Fine-Tuning.} Hard negatives---phonetically similar names from different places---sharpen discrimination. The loss is the same triplet form as Phase~1 but applied to Student embeddings:
\begin{equation}
\mathcal{L}_{\text{hard}} = \max\!\big(0,\; \|e_S^a - e_S^p\|_2 - \|e_S^a - e_S^n\|_2 + m\big)
\label{eq:phase3-loss}
\end{equation}
with $m = 0.3$. No distillation component is retained. Hard negatives are constructed with high orthographic similarity to anchors (same script, same 2-character prefix) but no shared place attestations. Training uses AdamW ($\text{lr} = 5 \times 10^{-5}$, batch size 1024) over 30 epochs (7.5h, final val\_loss 0.0212). The trade-off is deliberate: hard negatives from the same script teach finer within-script discrimination at modest cost to cross-script performance; same-script variants can be handled by traditional edit-distance methods.

Embedding inference for 67M toponyms required 2.5h. Final embeddings are quantised to \texttt{int8} and bulk-indexed to Elasticsearch. Total pipeline execution spans approximately 4 days wall-clock time.


\subsection{Training Data}
\label{subsec:data}

\subsubsection{Data Selection and Curation}
\label{subsec:selection}

Training data are extracted from three major gazetteer authorities---GeoNames, Wikidata, and the Getty TGN---via an existing consolidated index of 47.1 million place records. From these, 112.0 million toponym records spanning 1,944 languages and 20 script categories are extracted. After filtering 1.77 million pre-romanised forms and deduplication, \textbf{66.9 million unique toponyms} remain. The training subset from the three namespaces yields 57.6 million toponyms.

Four curation criteria address data quality: (1)~\textbf{Stratified sampling} by script+language pair (capped at 50,000 per bin, with oversampling up to 5$\times$ for small bins) prevents high-resource dominance; (2)~\textbf{Global vocabulary construction} scans the entire 66.9M corpus, yielding 113,280 tokens across 20 scripts; (3)~\textbf{Cross-script pairs} emerge naturally from density-based clustering (HDBSCAN~\citep{mcinnes2017hdbscan} with $\epsilon = 0.2$) on articulatory feature embeddings, generating positive pairs only from phonetically coherent groups within each place record; (4)~\textbf{Place-local deduplication} allows cross-place duplicates while preventing identical pairs within a single place's cluster.

The clustering approach correctly separates phonetically distinct name families. For example, a Cologne record yields two clusters: Germanic (Köln, Keulen) and Romance (Cologne, Colonia). Pairs are generated within clusters (Köln--Keulen, Cologne--Colonia) but not across them (Köln--Cologne), preventing the model from learning that phonetically unrelated exonyms should cluster together.

A \textbf{London} place record with 21 multilingual toponyms yields a single dominant cluster of 17 members spanning Arabic, CJK, Cyrillic, and Latin scripts (intra-cluster similarity 0.91): London (de, vi, pl, hu, cs, es, it, tr, nl, sv, en), \textar{لندن} (fa, ar, ur), Лондон (ru, uk), and \textcn{伦敦} (zh) cluster together, while French/Portuguese London (fr, pt), Bengali \textbn{লন্ডন}, and Serbian Ландон are correctly isolated as phonetically distinct variants. Similar multi-cluster patterns emerge for Moscow (2 major clusters), Beijing (3 clusters for ``Beijing'', ``Peking'', ``Pechino''), and Paris (3 clusters across 7 scripts). For places with only two toponyms, a cosine similarity threshold of 0.5 on PanPhon192 embeddings is used as fallback.

For homonym disambiguation, a candidate negative is valid only if it shares \emph{no place attestations} with the anchor, preventing names like ``Springfield'' from being used as negatives for other Springfields that may refer to the same location.

\subsubsection{IPA Transcription and Feature Extraction}
\label{subsec:g2p}

IPA transcriptions are generated via three complementary backends: Epitran~\citep{mortensen2018} (extended with 102 new language-script mappings), Phonikud for Hebrew, and CharsiuG2P~\citep{zhu2022charsiu-g2p} for Chinese topolects and Korean. Japanese Hiragana and Katakana are routed to Epitran by script detection prior to language-based routing.

\paragraph{Epitran Extensions.} Epitran ships with approximately 150 grapheme-to-phoneme map files covering perhaps 80--90 distinct language-script pairs. Since the corpus contains toponyms in over 1,900 language codes, 102 additional extension files were developed covering languages present in the corpus but lacking native Epitran support, spanning 14 scripts across diverse language families including Austronesian (Acehnese, Balinese, Sundanese), Celtic (Breton, Welsh, Irish), and Turkic (Bashkir, Chuvash, Crimean Tatar). For each target language-script pair, grapheme-to-phoneme rules were drafted by prompting multiple commercial large language models in rotation, cross-checking outputs for consistency and iterating until convergence. This process is best understood as knowledge distillation under noise: the Teacher-Student architecture is specifically designed to learn from imperfect training signals, since the Student learns to smooth over the Teacher's artifacts through the distillation and hard-negative phases. Individual grapheme-to-phoneme errors in extension files therefore propagate as stochastic noise rather than systematic bias in the final embeddings. Languages with extension-based G2P exhibit comparable intra-cluster cosine similarity distributions to those with native Epitran support during the HDBSCAN clustering stage (Section~\ref{subsec:selection}), providing indirect evidence that rule quality is sufficient for downstream training. Each extension file is a CSV mapping graphemes (or grapheme sequences) to IPA transcriptions, following Epitran's standard format. Intrinsic evaluation of G2P accuracy for each extension would require native-speaker phonological judgements across all 102 languages, which was beyond scope. Instead, the extensions are validated extrinsically: IPA transcriptions feed the Teacher network's PanPhon features, and the distillation architecture absorbs residual G2P errors as described above. Scripts relying heavily on extended Epitran achieve strong cross-script pass rates on systematically sampled pairs: Greek (93\% of tested pairs above 0.75 similarity), Armenian (94\%), and Gujarati (94\%), despite these scripts representing less than 1\% of training data (Table~\ref{tab:script-distribution}). However, as noted, individual G2P errors are tolerated rather than corrected by the architecture, and performance on languages where the extensions are least accurate may be correspondingly weaker. The extension files are released with the model to enable reproduction and community improvement.

PanPhon~\citep{mortensen2016} converts IPA segments into 24-dimensional articulatory feature vectors. To obtain fixed-length representations, 8-bin positional pooling divides each sequence into 8 bins and averages features within each, yielding $8 \times 24 = 192$ dimensions (PanPhon192). This preserves positional information while mapping variable-length sequences to fixed representations. Overall IPA coverage is 54.0\% (31.1M of 57.6M training-namespace toponyms); names without G2P are excluded from Teacher training but processed by the Student, which generalises from related scripts.

\subsubsection{Dataset Statistics}

The script distribution across all 66.9M toponyms and IPA coverage within the 57.6M training-namespace toponyms are shown in Table~\ref{tab:script-distribution}. From 8.2 million places with at least two toponyms, HDBSCAN clustering generates 65.1 million positive pairs across 595 script:language bins. After bin-balancing, 27.6 million pairs yield \~{}20.4 million Phase~1 training triplets and \~{}8 million Phase~3 hard-negative triplets.


\section{Results}
\label{sec:results}

\subsection{Embedding Quality}

Table~\ref{tab:embedding-quality} summarises embedding quality across diagnostic categories. The production index achieves 100\% embedding coverage over all 66.9M toponyms. Representative cross-script similarities include: London/Лондон 0.991, Athens/Αθήνα 0.980, Beijing/\textcn{北京} 0.955, Baghdad/\textar{بغداد} 0.969, Jerusalem/\texthebrew{ירושלים} 0.892. Cross-language same-script pairs with genuinely different pronunciations receive appropriately low scores (London/Londres 0.474, reflecting English /ˈlʌndən/ vs French /lɔ̃dʁ/), demonstrating correct phonetic discrimination. All four diacritic variant tests pass with similarities $\geq$0.95.

\subsection{MEHDIE Benchmark}

Symphonym is evaluated on the MEHDIE Hebrew-Arabic historical toponym benchmark~\citep{mehdie2025} using ranking metrics that reflect the system's design as a candidate generator. Table~\ref{tab:mehdie-ranking} presents comparative performance across all five testsets.

\paragraph{Ranking Performance.} The table includes four methods: PanPhon192 (raw articulatory features, no neural training), AnyAscii-augmented Levenshtein and Jaro-Winkler string baselines, and Symphonym. PanPhon192 serves as an ablation: it uses the same G2P pipeline and positional pooling that produce the Teacher's input features, but with raw cosine similarity only. This isolates the contribution of the training curriculum.

PanPhon192 achieves only 41.1\% R@1 and 45.0\% MRR---roughly half the trained system's performance. The string baselines substantially outperform PanPhon192 (Levenshtein: 81.5\% R@1, 88.5\% MRR), demonstrating that generic romanisation with edit distance is a stronger cross-script heuristic than raw phonetic features without learned alignment. Symphonym achieves the highest mean Recall@1 (85.2\%) and MRR (90.8\%), winning R@1 on three of five testsets with particularly large margins on TS9 (94.4\% vs Levenshtein 77.8\%) and TS10 (72.7\% vs 66.7\%). At Recall@10, Symphonym achieves 97.6\%.

\paragraph{Testset Analysis.} TS8 and TS9 yield near-perfect performance (R@1: 95.2--94.4\%, R@10: 100\%). TS10 (Yaqut-Kima Maghreb) proves most challenging for all methods; Symphonym's 72.7\% R@1 still leads Levenshtein (66.7\%) and Jaro-Winkler (54.5\%). The Maghreb toponyms contain more phonetically divergent historical variants reflecting genuine pronunciation evolution.

\paragraph{Comparison with MEHDIE System.} Direct comparison with the MEHDIE matching system~\citep{mehdie2025} is not straightforward: their evaluation uses threshold-based F-5 metrics (recall weighted 5$\times$ over precision), reflecting their users' preference for high recall at the cost of precision, which is incommensurable with ranking metrics. Their pipeline---a specialist Hebrew↔Arabic transliteration library with Phonetisaurus G2P and PanPhon Hamming distance---is carefully engineered for the phonetic kinship between these two Semitic languages. Symphonym addresses a different scope: general-purpose embedding across 20 scripts. The MEHDIE benchmark testsets---carefully verified ground-truth matches curated by domain experts---are the most directly relevant contribution from that work for evaluation purposes.

\subsection{Production Deployment}

The model was deployed to a staging Elasticsearch instance containing all 66,924,548 toponyms. To comprehensively evaluate cross-script matching capability, 11,723 genuine cross-script toponym pairs were sampled from the training data by systematically drawing up to 10 samples from each of 170+ cross-script script-pair bins (e.g., LATIN-CYRILLIC, ARABIC-BENGALI, CJK-HANGUL). These pairs derive from the same gazetteer sources used for training; this evaluation therefore tests whether trained embeddings produce correct similarity rankings when retrieved over the full 67M-toponym index, not whether the model generalises to unseen sources---the MEHDIE benchmark (Section~\ref{sec:results}, Table~\ref{tab:mehdie-ranking}) provides that independent evaluation on historical material absent from training. The 0.75 cosine similarity threshold is designed as a high-recall pruning threshold: it admits the vast majority of true cross-script equivalents into downstream processing while excluding only pairs with genuinely low phonetic correspondence---and testing yields 90.7\% pass rate (10,635/11,723 pairs exceed the threshold), with no missing documents (0 missing out of 66.9M indexed). The similarity distribution is strongly right-skewed: the majority of cross-script pairs achieve $\geq$0.90 similarity, with the densest concentration between 0.92 and 0.99. Pairs below 0.75 are concentrated in script combinations involving CJK-Hiragana (reflecting Mandarin vs Japanese phonetic mismatch) and the ``OTHER'' category (unsupported scripts). The CJK-Hiragana result (mean 0.437) reflects genuine phonological divergence rather than embedding instability: CharsiuG2P produces Mandarin phonetic readings for CJK characters, while Epitran produces Japanese readings (on'yomi or kun'yomi), and these are often phonetically unrelated for the same character. This is a phonological feature of the data, not a model deficiency. Importantly, even in lower-performing script-pair bins, the majority of sub-threshold pairs cluster near the 0.75 boundary, suggesting that downstream geographic filtering would resolve most ambiguous cases in a production pipeline.

\textbf{Performance by script pair:} Best-performing combinations include Hiragana-Katakana (mean 0.981), Devanagari-Kannada (0.976), Devanagari-Telugu (0.976), Cyrillic-Latin (0.923, $n$=1,334), and Arabic-Latin (0.898, $n$=800). More challenging combinations include CJK-Latin (0.808, $n$=1,073) and CJK-Hiragana (0.437, $n$=65), discussed above. Analysis of 10,000 randomly sampled embeddings confirms that the embedding space has desirable geometric properties for large-scale retrieval. Mean L2 norm is 1.000 $\pm$ 0.002, confirming that embeddings lie on the unit hypersphere as intended. Mean pairwise cosine similarity is 0.059, indicating that the space is neither collapsed (which would produce uniformly high similarity, degrading retrieval precision) nor excessively sparse. These properties are critical for operational Digital Earth systems: in high-dimensional approximate nearest-neighbour indices, embedding space degradation---whether through norm drift or dimensional collapse---produces cascading retrieval failures that are difficult to diagnose at scale.

\paragraph{KNN Retrieval.} Unlike pairwise similarity tests, KNN retrieval evaluates whether the embedding space naturally clusters variants together in open-ended searches across the full 67M corpus. Testing on representative queries shows robust cross-script retrieval: a ``London'' query retrieves Лондон (Cyrillic, 0.997), \textar{لندون} (Arabic, 0.991), and \textgeorgian{ლონდონი} (Georgian, 0.989), while correctly excluding phonetically distinct Romance variants (Londres, Londra). This exclusion reflects correct phonetic behaviour, not a deficiency: English /ˈlʌndən/ differs substantially from French /lɔ̃dʁ/ (nasal vowel, uvular fricative). Such phonetically divergent variants are linked through the complementary places index rather than phonetic similarity (see Discussion).

Production deployment must account for two practical challenges. First, high-multiplicity clusters: ``London'' appears in 69 language variants with near-identical embeddings, dominating top-$k$ results; post-processing strategies (script diversity re-ranking, candidate expansion with geographic filtering) are necessary. Second, length sensitivity: OpenStreetMap-derived institutional names (e.g., ``Kachua-mokampukur F P School'') can produce spurious high-similarity matches against short toponyms due to the PanPhon192 binning. The Student encoder's length bucket embedding mitigates this in practice, but the phenomenon illustrates the importance of length-aware post-filtering in production systems.

\paragraph{Scalability.} Symphonym embeddings integrate with Elasticsearch's HNSW approximate nearest-neighbour indexing~\citep{malkov2020hnsw}. Query latency averages 15--50ms over 67M toponyms. The Student model's inference cost is minimal: encoding a toponym requires a single forward pass through an 8.3M-parameter network ($<$1ms on CPU).

\section{Discussion}
\label{sec:discussion}

The evaluation results support Symphonym's central claim: a character-level neural encoder can learn phonetically meaningful representations across writing systems without requiring phonetic transcription at inference time.

\paragraph{Cross-script Generalisation.} The MEHDIE evaluation is particularly notable because the benchmark sources---medieval Hebrew and Arabic geographical texts---are independent of the modern gazetteers used for training. Strong cross-temporal performance suggests the model has learned general phonetic mappings rather than memorising specific toponyms. This generalisation emerges from the Teacher-Student architecture: the Teacher learns language-specific phonetic mappings from IPA features, then transfers this knowledge through distillation. The result has direct implications for geospatial data integration: a model trained on modern authority data generalises to historical sources without specialist tuning, addressing a long-standing barrier to integrating archival geographic material with contemporary databases.

\paragraph{Value of Neural Training.} The PanPhon192 ablation (Table~\ref{tab:mehdie-ranking}) confirms that Symphonym's performance derives from the three-phase training curriculum rather than the articulatory features alone. PanPhon192's mean MRR of 45.0\% is less than half Symphonym's 90.8\%, and substantially below even the string baselines (Levenshtein 88.5\%). The gap is largest on TS7 and TS10 (MRR 19.2\% and 16.6\%), where raw articulatory features cannot bridge the phonetic distance between medieval Arabic and Hebrew variants without learned alignment. Even on TS9, where PanPhon192 performs best (82.0\% MRR), Symphonym still improves by 15 points (97.2\%). This result has methodological implications: articulatory features provide a necessary phonetic grounding---without them, the model would have no cross-script signal to learn from---but they are radically insufficient on their own. The neural architecture learns a non-linear alignment between articulatory feature spaces across languages that simple distance metrics on raw features cannot capture.

\paragraph{Same-script Performance.} Same-script cross-language pairs achieve 86\% pass rate, reflecting correct phonetic discrimination: exonyms with genuinely different pronunciations receive appropriately low scores. This behaviour is operationally desirable---a phonetic index \emph{should not} link ``Germany'' to ``Deutschland'' or ``\textcn{東京}'' to ``\textjp{とうきょう}'', because these are phonetically unrelated despite referring to the same place. In the production system, such links are provided by a separate \emph{places} index that groups all toponyms attested for the same geographic entity regardless of phonetic similarity. The places index thus complements the phonetic embeddings: Symphonym handles cross-script matching where names \emph{sound alike} (Baghdad/\textar{بغداد}/Багдад); the places index handles cases where names \emph{refer to the same place} but sound different (Germany/Deutschland, or CJK Kanji via Mandarin phonetics vs Hiragana via Japanese readings). Together, the two indices cover the full range of cross-script and cross-language name variation encountered in gazetteer integration. Traditional string methods (Levenshtein, Jaro-Winkler) provide a third complementary channel for same-script refinement.

\paragraph{Operational Implications.} In a planetary-scale SDI, comparing a query against every indexed record is infeasible. Symphonym's primary operational contribution is search space reduction: approximate nearest-neighbour retrieval over the HNSW index reduces 67 million candidates to a small set of phonetically plausible matches in under 50ms, after which expensive spatial and temporal disambiguation can be applied. This layered retrieval architecture---phonetic prior, then geographic filter---preserves high recall while bounding computational cost. Unlike rule-based phonetic algorithms (Soundex, Metaphone), which require exact code matches and operate within single scripts, embedding similarity enables fuzzy matching with graceful degradation across all scripts simultaneously. The Student encoder requires only raw character input---no language identification, no G2P conversion, no phonetic databases---which simplifies deployment and eliminates failure modes associated with unknown languages or unsupported scripts. Within a production reconciliation pipeline, cross-script candidate retrieval uses phonetic embeddings; same-script refinement may use edit distance; and final disambiguation incorporates geographic proximity and temporal constraints. The embedding similarity score thus provides a phonetic prior that downstream components weigh against other evidence.

\subsection{Limitations}
\label{subsec:limitations}

\paragraph{Training Data Coverage.} Despite stratified sampling, training sources retain geographic biases. GeoNames over-represents populated places with official names; Wikidata skews toward places of encyclopaedic interest; TGN emphasises art-historically significant locations. Performance on under-represented scripts and mundane places lacking multilingual attestations may be weaker.

\paragraph{Tonal Languages.} The architecture does not explicitly model tone, which is phonemically contrastive in Chinese, Vietnamese, and Thai. PanPhon encodes segmental articulatory properties but not suprasegmental features. In practice, tonal minimal pairs are rare in geographic naming, and Chinese place names have distinct segmental content that Symphonym captures effectively (0.97--0.99 similarity with romanised forms).

\paragraph{Confusable Pairs.} Phonetically similar strings receive high similarity regardless of semantic relationship (Austria/Australia: 0.883, China/Ghana: 0.932). Disambiguation requires geographic or contextual evidence beyond phonetic similarity.

\subsection{Cross-Temporal Generalisation and Broader Applicability}

A comprehensive Digital Earth must integrate data across the temporal as well as the spatial axis. The MEHDIE benchmark results (Section~\ref{sec:results}) already demonstrate cross-temporal transfer from modern training data to medieval sources; the following examples extend this to pre-standardisation orthographic variation, suggesting that the model functions as an infrastructure tool for long-duration geospatial knowledge graphs.

Beyond cross-script matching, Symphonym naturally handles pre-standardisation spelling variation---the inconsistent orthography characteristic of historical geographic sources. Colonial-era survey records, historical mapping projects, and archival catalogues routinely exhibit this variation, and phonetically grounded embeddings cluster sound-alike variants near their modern canonical forms without language-specific spelling normalisation rules.

A case study on medieval London merchant names demonstrates successful clustering of phonetically similar spelling variants (e.g., multiple permutations of ``Deryke/Derico/Diryk Shotynbaker/Shuttynbaker/Shotyngbaker'') without retraining, suggesting applicability beyond toponyms to any geospatial metadata integration task where name forms vary across archival sources. In geospatial contexts, such name resolution enables linking individuals to places across archival sources—mapping historical trade networks, property transfers, or migration patterns where the same person appears under orthographically variable name forms in different records.

By learning phonetic correspondences that are difficult to enumerate manually, Symphonym's fixed-length embeddings address a practical and growing challenge in geospatial data integration: linking name references across languages, scripts, and historical periods at production scale. The approach generalises beyond toponyms to other named entity classes encountered in SDI contexts---hydronyms, ethnonyms, and institutional names---and to linked open data reconciliation tasks where URI-based linkage is unavailable because matching records have not yet been identified.

\subsection*{Acknowledgments}

This research used the HTC and H2P clusters at the University of Pittsburgh Center for Research Computing and Data (RRID:SCR\_022735), supported by NIH award S10OD028483 and NSF award OAC-2117681 respectively.

In accordance with Taylor \& Francis policy, the following uses of generative AI tools are disclosed. Claude Sonnet 4.6 (Anthropic), GPT-5 (OpenAI), and Gemini 1.5 Pro (Google) were used to draft grapheme-to-phoneme extension files for the Epitran library (Section~2.5.2), with outputs cross-checked for consistency and validated extrinsically through downstream performance. Claude and GitHub Copilot (Microsoft) were used for coding assistance during system development, and Claude for structural feedback and language editing during manuscript preparation. All content was reviewed, validated, and revised by the author, who takes full responsibility for the accuracy and integrity of the work.

\subsection*{Declaration of Interest Statement}

No relevant financial or non-financial interests to disclose.

\subsection*{Data Availability Statement}

All trained models, vocabularies, extension files, and evaluation results are openly available at \url{https://doi.org/10.5281/zenodo.18682017} \citep{symphonym2026zenodo}. Training code and inference utilities are available at \url{https://huggingface.co/docuracy/symphonym-v7} \citep{symphonym2026hf} and \url{https://github.com/WorldHistoricalGazetteer/whgazetteer}. An expanded preprint covering additional architectural detail and analysis of the training pipeline is available as~\citet{symphonym2026}. Training data are derived from publicly available sources: GeoNames (\url{https://www.geonames.org/}), Wikidata (\url{https://www.wikidata.org/}), and Getty TGN (\url{https://www.getty.edu/research/tools/vocabularies/tgn/}). The MEHDIE benchmark is available via \citet{mehdie2025}.


\bibliographystyle{plainnat}
\bibliography{references}


%% ============================================================
%% TABLES — each on a separate page, after references
%% ============================================================

\clearpage

\begin{table}[p]
\centering
\small
\begin{tabular}{lrrrl}
\toprule
\textbf{Script} & \textbf{Count} & \textbf{\%} & \textbf{IPA Coverage} & \textbf{Top Languages (IPA)} \\
\midrule
LATIN & 55,617,677 & 83.1\% & 49.8\% & en, fr, nl, de, sv \\
CYRILLIC & 3,614,762 & 5.4\% & 47.1\% & ru, uk, bg, sr \\
CJK & 2,973,525 & 4.4\% & 50.1\% & zh (CharsiuG2P) \\
ARABIC & 2,098,089 & 3.1\% & 52.5\% & fa, ar, ur \\
HANGUL & 393,996 & 0.6\% & 58.0\% & ko (CharsiuG2P) \\
OTHER & 342,642 & 0.5\% & 0.0\% & --- \\
KATAKANA & 340,555 & 0.5\% & 91.2\% & ja (\texttt{jpn-Kana}) \\
THAI & 251,458 & 0.4\% & 83.6\% & th \\
GREEK & 217,997 & 0.3\% & 77.4\% & el (Epitran ext.) \\
DEVANAGARI & 166,957 & 0.2\% & 57.2\% & hi, mr, ne \\
ARMENIAN & 153,467 & 0.2\% & 93.7\% & hy (Epitran ext.) \\
HIRAGANA & 151,980 & 0.2\% & 31.3\% & ja (\texttt{jpn-Hira}) \\
HEBREW & 151,960 & 0.2\% & 83.8\% & he (Phonikud) \\
BENGALI & 106,896 & 0.2\% & 72.9\% & bn \\
GEORGIAN & 105,902 & 0.2\% & 81.2\% & ka \\
MALAYALAM & 68,176 & 0.1\% & 78.5\% & ml \\
TAMIL & 52,486 & 0.1\% & 90.9\% & ta \\
TELUGU & 51,440 & 0.1\% & 92.6\% & te \\
KANNADA & 43,155 & 0.1\% & 48.6\% & kn (Epitran ext.) \\
GUJARATI & 21,428 & 0.03\% & 94.9\% & gu (Epitran ext.) \\
\bottomrule
\end{tabular}
\caption{Script distribution across all 66.9M unique toponyms. IPA coverage percentages reflect successful transcription within the 57.6M training namespace toponyms (gn/wd/tgn). Overall IPA coverage is 54.0\% (31.1M toponyms).}
\label{tab:script-distribution}
\end{table}

\clearpage

\begin{table}[p]
\centering
\small
\begin{tabular}{p{3.2cm} c p{5.8cm}}
\toprule
\textbf{Test Category} & \textbf{Pass Rate} & \textbf{Description} \\
\midrule
Cross-script equivalents & 18/22 (81.8\%) & Latin $\leftrightarrow$ Cyrillic, Greek, Arabic, CJK, Hebrew, Hangul \\
Diacritic variants & 4/4 (100\%) & Zurich/Zürich, Krakow/Kraków, São Paulo/Sao Paulo \\
Unrelated pairs & 3/3 (100\%) & Correctly separated (low similarity) \\
\midrule
\textbf{Total} & \textbf{25/29 (86.2\%)} & \\
\bottomrule
\end{tabular}
\caption{Production embedding quality diagnostics. Cross-script matching is the primary design goal. Tested on 66.9M toponyms with 100\% embedding coverage across 20 scripts.}
\label{tab:embedding-quality}
\end{table}

\clearpage

\begin{table}[p]
\centering
\small
\begin{tabular}{llrrrr}
\toprule
\textbf{Method} & \textbf{Testset} & \textbf{R@1} & \textbf{R@5} & \textbf{R@10} & \textbf{MRR} \\
\midrule
\multirow{6}{*}{\textbf{PanPhon192}}
& TS7 (Yaqut-Kima Sham) & 15.2 & 21.2 & 27.3 & 19.2 \\
& TS8 (Kima-Thurayya Sham) & 28.6 & 28.6 & 42.9 & 32.3 \\
& TS9 (Tudela-Thurayya) & 77.8 & 88.9 & 88.9 & 82.0 \\
& TS10 (Yaqut-Kima Maghreb) & 12.1 & 21.2 & 21.2 & 16.6 \\
& TS11 (Damast-Tudela) & 71.9 & 81.2 & 81.2 & 75.1 \\
\cmidrule{2-6}
& \textbf{Mean} & \textbf{41.1} & \textbf{48.2} & \textbf{52.3} & \textbf{45.0} \\
\midrule
\multirow{6}{*}{\textbf{Levenshtein}}
& TS7 (Yaqut-Kima Sham) & 75.8 & 93.9 & 100.0 & 83.7 \\
& TS8 (Kima-Thurayya Sham) & 90.5 & 100.0 & 100.0 & 94.4 \\
& TS9 (Tudela-Thurayya) & 77.8 & 100.0 & 100.0 & 87.5 \\
& TS10 (Yaqut-Kima Maghreb) & 66.7 & 97.0 & 97.0 & 79.6 \\
& TS11 (Damast-Tudela) & 96.9 & 96.9 & 100.0 & 97.3 \\
\cmidrule{2-6}
& \textbf{Mean} & \textbf{81.5} & \textbf{97.5} & \textbf{99.4} & \textbf{88.5} \\
\midrule
\multirow{6}{*}{\textbf{Jaro-Winkler}}
& TS7 (Yaqut-Kima Sham) & 75.8 & 100.0 & 100.0 & 86.4 \\
& TS8 (Kima-Thurayya Sham) & 90.5 & 90.5 & 95.2 & 91.4 \\
& TS9 (Tudela-Thurayya) & 77.8 & 100.0 & 100.0 & 86.6 \\
& TS10 (Yaqut-Kima Maghreb) & 54.5 & 93.9 & 97.0 & 71.9 \\
& TS11 (Damast-Tudela) & 93.8 & 96.9 & 96.9 & 95.4 \\
\cmidrule{2-6}
& \textbf{Mean} & \textbf{78.5} & \textbf{96.2} & \textbf{97.8} & \textbf{86.3} \\
\midrule
\multirow{6}{*}{\textbf{Symphonym}}
& TS7 (Yaqut-Kima Sham) & 69.7 & 97.0 & 100.0 & 82.9 \\
& TS8 (Kima-Thurayya Sham) & 95.2 & 100.0 & 100.0 & 96.8 \\
& TS9 (Tudela-Thurayya) & 94.4 & 100.0 & 100.0 & 97.2 \\
& TS10 (Yaqut-Kima Maghreb) & 72.7 & 87.9 & 87.9 & 80.8 \\
& TS11 (Damast-Tudela) & 93.8 & 100.0 & 100.0 & 96.4 \\
\cmidrule{2-6}
& \textbf{Mean} & \textbf{85.2} & \textbf{97.0} & \textbf{97.6} & \textbf{90.8} \\
\bottomrule
\end{tabular}
\caption{MEHDIE benchmark ranking metrics (all values in \%). PanPhon192 is the raw 192-dimensional articulatory feature representation used as \emph{input} to the Teacher---phonetic features before any neural training. String baselines are augmented with AnyAscii romanisation. Symphonym achieves the highest mean R@1 (85.2\%) and MRR (90.8\%), doubling PanPhon192's MRR (45.0\%) and outperforming both string baselines.}
\label{tab:mehdie-ranking}
\end{table}


\end{document}

